\section{Introdução e Objetivos}
    \subsection{Introdução}
    A termodinâmica química estuda as trocas de energia que acompanham as transformações da matéria. Pela Primeira Lei da Termodinâmica, a energia se conserva, de modo que o calor liberado ou absorvido por uma reação é transferido ao ambiente ao seu redor\textsuperscript{[1]}. A grandeza que descreve o calor trocado a pressão constante é a variação de entalpia ($\Delta H$): reações com $\Delta H < 0$ liberam calor (exotérmicas) e reações com $\Delta H > 0$ absorvem calor (endotérmicas)\textsuperscript{[2]}.

    A calorimetria é a técnica que mede esse calor a partir da variação de temperatura de uma massa conhecida de solução em um calorímetro. Em um calorímetro ideal, todo o calor da reação seria absorvido pela solução, o que permite escrever $\Delta H_{\text{reação}} = -\Delta H_{\text{solução}} = -m_s\,C_s\,\Delta T$, em que $m_s$ é a massa da solução, $C_s$ seu calor específico e $\Delta T$ a variação de temperatura\textsuperscript{[1]}. Como a entalpia é uma função de estado, o calor de uma reação que ocorre em etapas é a soma dos calores das etapas individuais, resultado conhecido como Lei de Hess\textsuperscript{[2]}.

    Neste experimento, a dissolução do hidróxido de sódio sólido, a reação do $\mathrm{NaOH}$ sólido com ácido clorídrico e a neutralização entre soluções aquosas de $\mathrm{HCl}$ e $\mathrm{NaOH}$ foram estudadas em um calorímetro de poliestireno, de modo a relacionar os calores dessas reações por meio da Lei de Hess.

    \subsection{Objetivos}
    Os objetivos deste experimento foram determinar, por calorimetria, as entalpias das reações (A) dissolução do $\mathrm{NaOH}$ sólido em água, (B) reação entre $\mathrm{HCl}$ aquoso e $\mathrm{NaOH}$ sólido e (C) neutralização entre $\mathrm{HCl}$ aquoso e $\mathrm{NaOH}$ aquoso, e verificar a relação entre elas por meio da Lei de Hess, discutindo as principais fontes de erro do método.
